\section{Future Directions}

\subsection{Protocol Boundary}

The present paper stops at the mathematical and formalization boundary. A
subsequent protocol paper should turn the commitment, oracle, proof-carrying,
identity, storage, and proximity surfaces into explicit adversary games and
executable verification paths. The concrete protocol must specify which
certificates are public evidence, which witnesses remain private, and which
runtime failures cause rejection, hold, or re-evaluation.

On the identity side, the hidden-representative/public-classifier split should
be developed into explicit class-membership, unlinkability, and nullifier
properties. Public identity data should remain class values, hashes, proofs, or
oracle judgments; the private side should retain fibration representatives,
transport data, and the Hodge operator selecting the operative kernel.

On the storage side, the commitment layer can be developed as a controlled
system for sheaf/cochain projections. A public store may retain topology
hashes, cochain roots, compressed public data, and proof-carrying statements, while
full sheaf and fibration witnesses remain private or are opened only through
dedicated proof paths.

The protocol-facing reference point is the higher categorical cryptoeconomics
program \cite{higherCategoricalCryptoeconomics}. It is relevant to the design
of protocol evidence and state semantics, but it is not used in the proofs of
the present paper.

\subsection{Complexity-Theoretic Refinements}

The current formalization isolates the hidden-operator obstruction: public data
does not canonically determine the operative Hodge--Laplacian kernel. It also
separates the external known-lattice CVP assumption from the project-native
hidden-operator ambiguity, gives public-data-only oracle insufficiency
statements, records classical and quantum adversary likelihood surfaces, and
proves infinitude results for the relevant Grassmannian/Laplacian families.
Future work should sharpen these surfaces quantitatively: rates of growth for
hidden-kernel families, refined leakage bounds for public classifiers, success
probabilities for specified adversary classes, and reductions between variants
of HCP.
